% ===========================================================
%  第 6 章 向量组与向量空间 —— 高等代数【深化】拓展框（转录补充）
%  来源：北大《高等代数》第五版（偏移 PDF = 印刷 + 12）
%        樊启斌《高等代数典型问题与方法》（偏移 PDF = 印刷 + 9）
%  说明：主控已在 deep_ch06_authored.tex 写过 3 块
%        （过渡矩阵与坐标变换 / 施密特正交化 / 正交矩阵），
%        同章另有 deep_ch06_fan.tex（樊 §9.2 正交矩阵与正交变换、§9.4 Gram 矩阵），
%        本文件只补充它们未覆盖的面向，不重复已有条目。
%        （度量矩阵一块已改写为“坐标表达式 + 换基合同”，
%         其定义与正定性交给 deep_ch06_fan.tex 的 Gram 矩阵块，避免撞车。）
%  对应考纲〔向量〕第 5/6/7/8 条（9 讲覆盖薄弱、标记"必须补强"）
% ===========================================================

% ---------- 线性空间的定义（北大 六.2 印刷 163–165；樊 5.1.1 印刷 272）----------
\begin{fdeep}{线性空间：从具体对象到八条公理}
设 $V$ 是一个非空集合，$P$ 是一个数域。在 $V$ 的元素之间定义了加法，在 $P$ 与 $V$ 之间定义了数量乘法。若这两种运算满足下面八条规则，就称 $V$ 是数域 $P$ 上的\textbf{线性空间}（也叫向量空间），$V$ 的元素统称\textbf{向量}。

\textbf{加法四条}：$\alpha+\beta=\beta+\alpha$；$(\alpha+\beta)+\gamma=\alpha+(\beta+\gamma)$；$V$ 中存在零元素 $\mathbf{0}$，使 $\mathbf{0}+\alpha=\alpha$；每个 $\alpha$ 存在负元素 $-\alpha$，使 $\alpha+(-\alpha)=\mathbf{0}$。

\textbf{数乘与分配四条}：$1\alpha=\alpha$；$k(l\alpha)=(kl)\alpha$；$(k+l)\alpha=k\alpha+l\alpha$；$k(\alpha+\beta)=k\alpha+k\beta$。

\textbf{公理化的意义}：$P$ 上全体 $n$ 元有序数组、全体 $m\times n$ 矩阵、次数小于 $n$ 的多项式、闭区间上的连续函数，都满足这八条，于是关于线性空间证明的每一条结论可以\textbf{一次性地}适用于全部这些对象。

\textbf{三条可直接推出的简单性质}：零元素唯一；负元素唯一；$0\alpha=\mathbf{0}$，$k\mathbf{0}=\mathbf{0}$，$(-1)\alpha=-\alpha$；并且若 $k\alpha=\mathbf{0}$，则必有 $k=0$ 或 $\alpha=\mathbf{0}$。
【反哺点】服务考纲〔向量〕“理解 $n$ 维向量空间的概念”。它解释了考研中“向量”为什么不必是箭头：凡是带这两种运算且满足八条规则的对象都是向量，所以矩阵、多项式、函数都能当向量处理，矩阵的秩、方程组的解才能用“维数、基、坐标”统一叙述。判断某集合能否构成线性空间，只看加法与数乘是否封闭、零元与负元是否在集合内。
\end{fdeep}

% ---------- 子空间判别（北大 六.5 定义 7、定理 2，印刷 172–173）----------
\begin{fdeep}{子空间的判别：两条封闭性就够}
数域 $P$ 上线性空间 $V$ 的一个\textbf{非空子集} $W$ 称为 $V$ 的\textbf{线性子空间}（简称子空间），如果 $W$ 对 $V$ 的两种运算也构成数域 $P$ 上的线性空间。

由于八条公理中其余六条是从 $V$ 原样继承的，判别时只需验证两条封闭性：
\fml{\alpha,\beta\in W\ \Longrightarrow\ \alpha+\beta\in W,\qquad k\in P,\ \alpha\in W\ \Longrightarrow\ k\alpha\in W}
等价地可合并成一条：$W$ 非空，且对任意 $k,l\in P$、$\alpha,\beta\in W$ 都有 $k\alpha+l\beta\in W$。又因为取 $k=l=0$ 即得 $\mathbf{0}\in W$，实际做题时\textbf{先看零向量在不在 $W$ 里}——不含零向量的子集一定不是子空间。

\textbf{典型子空间}：只含零向量的零子空间 $\{\mathbf{0}\}$ 与 $V$ 本身（合称平凡子空间）；齐次线性方程组的全部解组成的\textbf{解空间}；由向量组生成的子空间 $L(\alpha_1,\dots,\alpha_s)$。任何子空间的维数都不超过整个空间的维数。
【反哺点】服务考纲〔向量〕“了解子空间的概念”。选择题常直接问“下列集合是否构成子空间”：先用“$\mathbf{0}$ 是否属于该集合”快速排除，再用一条 $k\alpha+l\beta$ 验证。高频陷阱是“平面上不过原点的直线”“满足 $|x|=1$ 的向量集”这类不含零向量的集合。
\end{fdeep}

% ---------- 生成子空间与秩（北大 六.5 定理 3，印刷 173–174；樊 5.1.4 印刷 273）----------
\begin{fdeep}{生成子空间 $L(\alpha_1,\dots,\alpha_s)$：极大无关组就是它的基}
向量组 $\alpha_1,\dots,\alpha_s$ 的所有线性组合组成的集合
\fml{L(\alpha_1,\dots,\alpha_s)=\{k_1\alpha_1+k_2\alpha_2+\cdots+k_s\alpha_s\mid k_i\in P\}}
是 $V$ 的一个子空间，称为由这组向量\textbf{生成}的子空间；它同时也是 $V$ 中一切包含 $\alpha_1,\dots,\alpha_s$ 的子空间里最小的那个。

两条关键结论（北大 六.5 定理 3）：
（1）$L(\alpha_1,\dots,\alpha_s)=L(\beta_1,\dots,\beta_t)$ 的充分必要条件是这两个向量组\textbf{等价}（能互相线性表出）；
（2）$\dim L(\alpha_1,\dots,\alpha_s)$ 等于向量组 $\alpha_1,\dots,\alpha_s$ 的\textbf{秩}。

第（2）条的证明只有一句话：取该向量组的一个极大线性无关组，它既线性无关，又能表出原组每个向量，因而是 $L$ 的一组基，而极大无关组的向量个数正是秩。
【反哺点】服务考纲〔向量〕“理解向量组的秩、极大线性无关组的概念”“了解子空间、基底、维数”。这一条把“向量组的秩”“极大线性无关组”“空间的基与维数”三个概念焊成一体：求 $L$ 的维数与基，就是求向量组的秩与一个极大无关组——把向量按列拼成矩阵、初等行变换化阶梯形，主元所在列对应的\textbf{原}向量组就是基，主元个数就是维数。
\end{fdeep}

% ---------- 线性相关/无关的定义与性质（北大 三.3 定义 11、11'、12，印刷 79–80）----------
\begin{fdeep}{线性相关与线性无关：两种定义与“部分、整体”性质}
\textbf{定义一}（存在表出）：向量组 $\alpha_1,\dots,\alpha_s$（$s\ge2$）中若有一个向量可由其余向量线性表出，就称它\textbf{线性相关}，否则称线性无关。

\textbf{定义二}（存在非零系数）：若数域 $P$ 中存在\textbf{不全为零}的数 $k_1,\dots,k_s$ 使
\fml{k_1\alpha_1+k_2\alpha_2+\cdots+k_s\alpha_s=\mathbf{0}}
则称线性相关；若上式只可能 $k_1=\cdots=k_s=0$，则称线性无关。

两个定义在 $s\ge2$ 时彼此等价：由 $\alpha_s=k_1\alpha_1+\cdots+k_{s-1}\alpha_{s-1}$ 立即得 $k_1\alpha_1+\cdots+k_{s-1}\alpha_{s-1}+(-1)\alpha_s=\mathbf{0}$，而 $-1\ne0$，系数不全为零；反过来，把非零系数 $k_s\ne0$ 的那一项解出来，就得到“可由其余向量表出”。

\textbf{四条常用性质}：1）部分相关 $\Longrightarrow$ 整体相关，等价的逆否说法是“整体无关 $\Longrightarrow$ 任何非空部分组无关”；2）含零向量的向量组必线性相关（零系数之外的系数可取零，故按定义一必相关）；3）单个向量 $\alpha$ 线性相关 $\iff$ $\alpha=\mathbf{0}$；4）线性无关组中不含两个成比例的向量。
【反哺点】服务考纲〔向量〕“理解向量组线性相关、线性无关的概念”。性质 1 与 2 是选择题的送分点：选项里出现“部分无关 $\Rightarrow$ 整体无关”或“整体相关 $\Rightarrow$ 部分相关”都是错的。另外要记住“成比例的向量必相关”，用于快速判断含重复列、含零列的矩阵不可逆。
\end{fdeep}

% ---------- 判定网（北大 三.3 印刷 81–82）----------
\begin{fdeep}{线性相关性的判定网：五条入口殊途同归}
设 $\alpha_1,\dots,\alpha_s$ 是 $n$ 维向量，把它们按列拼成 $n\times s$ 矩阵 $A=(\alpha_1,\dots,\alpha_s)$，则以下说法\textbf{两两等价}：
1）$\alpha_1,\dots,\alpha_s$ 线性相关；
2）齐次线性方程组 $x_1\alpha_1+x_2\alpha_2+\cdots+x_s\alpha_s=\mathbf{0}$ 有\textbf{非零解}；
3）$\operatorname{rank}A<s$；
4）其中至少有一个向量可由其余向量线性表出；
5）它的极大线性无关组所含向量个数小于 $s$。
当向量个数等于维数（$s=n$）时再添一条：$|A|=0$。此时“线性无关 $\iff$ $|A|\ne0$ $\iff$ $A$ 可逆”，三个说法可以随时互换。

另有两条只比“个数”的结论：$n+1$ 个 $n$ 维向量必线性相关；若 $\alpha_1,\dots,\alpha_r$ 可由 $\beta_1,\dots,\beta_s$ 线性表出，而 $\alpha_1,\dots,\alpha_r$ 线性无关，则必有 $r\le s$——这就是“以少表多必相关”。
【反哺点】服务考纲〔向量〕“掌握向量组线性相关、线性无关的有关性质及判别法”。这张网把行列式、秩、方程组、子空间维数四条路打通：考场上不必挑“正确”的一条，哪条入口给的数最少就用哪条。个数型（$s$ 与 $n$ 的大小比较）的题目直接套两条推论，是选择题最快的解法。
\end{fdeep}

% ---------- 极大无关组与秩（北大 三.3 定义 13、14、定理 3，印刷 83）----------
\begin{fdeep}{极大线性无关组与秩：它和“基”是同一个东西}
向量组的一个部分组若本身线性无关，并且从原组中再添进任何一个向量都线性相关，就称为原向量组的一个\textbf{极大线性无关组}；极大线性无关组所含向量的个数称为该向量组的\textbf{秩}。

三条要点：1）极大线性无关组必与原向量组\textbf{等价}，因为原组每个向量若不能由它表出，就可以添进去而仍然线性无关，与“极大”矛盾；2）一个向量组的\textbf{任意两个}极大线性无关组含相同个数的向量，所以“秩”是向量组自身的属性，与挑法无关；3）向量组线性无关 $\iff$ 它的秩等于它所含向量的个数。

\textbf{与“基”的打通}：向量组 $\alpha_1,\dots,\alpha_s$ 的一个极大线性无关组，恰好就是生成子空间 $L(\alpha_1,\dots,\alpha_s)$ 的一组\textbf{基}。于是“向量组的秩 = 它张成空间的维数”，第 3 章的“秩”与第 6 章的“维数”从此是同一件事的两个名字。
【反哺点】服务考纲〔向量〕“会求向量组的秩和极大线性无关组”。求法要点：把向量\textbf{按列}拼成矩阵，只作初等\textbf{行}变换化阶梯形；主元个数即秩，主元所在列对应的\textbf{原向量}构成极大无关组（不能用化简后的列向量作答）。反过来，若题目给出“秩为 $r$”，就等于说“存在 $r$ 个线性无关的向量且任意 $r+1$ 个相关”。
\end{fdeep}

% ---------- 向量组等价 vs 矩阵等价（北大 三.3 定义 10、印刷 78–79；六.5 定理 3）----------
\begin{fdeep}{向量组等价与矩阵等价：两个“等价”不是一回事}
\textbf{向量组等价}：两个向量组能\textbf{互相线性表出}。它满足自反性、对称性、传递性；并且向量组 $\alpha_1,\dots,\alpha_r$ 与 $\beta_1,\dots,\beta_t$ 等价当且仅当
\fml{L(\alpha_1,\dots,\alpha_r)=L(\beta_1,\dots,\beta_t)}
即等价与“生成同一个子空间”是同一件事。两条常用结论：等价的两个向量组秩必相等，但\textbf{秩相等推不出等价}（例如 $(1,0,0),(0,1,0)$ 与 $(1,0,0),(0,0,1)$ 的秩都是 $2$，张成的却是不同的平面）；两个都线性无关的等价向量组必含相同个数的向量。

\textbf{矩阵等价}：$A$ 经有限次初等变换化为 $B$，等价条件是 $A$ 与 $B$ \textbf{同型}（行数、列数分别相同）且秩相等。

\textbf{区别所在}：矩阵等价是同型矩阵之间的关系，判据只看秩；向量组等价是两组向量之间的关系，除了秩相等，还要求两个矩阵的列向量组张成\textbf{同一个}子空间。
【反哺点】服务考纲〔向量〕“理解向量组等价的概念”。判断题里“$A$ 与 $B$ 等价”与“$A$ 的列向量组与 $B$ 的列向量组等价”是两个不同的命题：前者只要同型同秩，后者要求列空间相同，因此“列向量组等价 $\Longrightarrow$ 矩阵等价”成立，反之不成立（除非 $A$ 可逆这类额外条件）。抓住这一点可避开一大类概念陷阱。
\end{fdeep}

% ---------- 维数唯一与扩充定理（北大 六.3 定理 1 印刷 167；六.5 定理 4 印刷 174）----------
\begin{fdeep}{有限维空间的基与维数：唯一性、判定与扩充}
\textbf{维数是空间自身的属性}。一个向量组的极大无关组可以不唯一，但所含向量个数相同；而 $n$ 维空间的两组基都线性无关且互相等价，由“线性无关的等价向量组含相同个数的向量”即知两组基的向量个数必然相等。所以 $\dim V$ 与基的选取无关。

\textbf{判定一组向量是否为基}（北大 六.3 定理 1）：若在线性空间 $V$ 中有 $n$ 个线性无关的向量 $\alpha_1,\dots,\alpha_n$，且 $V$ 中任一向量都能由它们线性表出，则 $V$ 是 $n$ 维的，且 $\alpha_1,\dots,\alpha_n$ 就是 $V$ 的一组基。证明思路：任取 $V$ 中 $n+1$ 个向量，它们都可由这 $n$ 个向量表出，由“以少表多必相关”知必线性相关，故 $V$ 中不存在 $n+1$ 个线性无关的向量。

\textbf{基的扩充定理}（北大 六.5 定理 4）：设 $W$ 是 $n$ 维空间 $V$ 的 $m$ 维子空间，$\alpha_1,\dots,\alpha_m$ 是 $W$ 的一组基，则在 $V$ 中必可找到 $n-m$ 个向量，使 $\alpha_1,\dots,\alpha_n$ 构成 $V$ 的一组基。
【反哺点】服务考纲〔向量〕“了解基底、维数的概念”。判定“某组向量是否为基”时，定理 1 把工作量减半：只需验“线性无关 + 能表出”，不必再验“$V$ 中每个向量表法唯一”。扩充定理则回答了常见追问“已给的无关组能不能补成基”——能，且补的个数恰为 $\dim V-\dim W$。
\end{fdeep}

% ---------- 子空间的交与和（北大 六.6 定义 8、定理 5、6，印刷 174–175）----------
\begin{fdeep}{子空间的交与和：两种运算与包含关系}
设 $V_1,V_2$ 是线性空间 $V$ 的子空间。
\textbf{交}：$V_1\cap V_2$ 仍是 $V$ 的子空间（验证：若 $\alpha,\beta\in V_1\cap V_2$，则 $\alpha,\beta$ 同属于 $V_1$ 也同属于 $V_2$，故 $k\alpha+l\beta$ 同属于两者）。多个子空间的交 $\bigcap_{i=1}^{s}V_i$ 也是子空间。
\textbf{和}：所有能表示成 $\alpha_1+\alpha_2$（$\alpha_1\in V_1$，$\alpha_2\in V_2$）的向量组成的集合称为 $V_1$ 与 $V_2$ 的\textbf{和}，记作 $V_1+V_2$，它也是 $V$ 的子空间。多个子空间的和 $\sum_{i=1}^{s}V_i$ 由所有形如 $\alpha_1+\cdots+\alpha_s$（$\alpha_i\in V_i$）的向量组成。

\textbf{运算律}：交与和都满足交换律与结合律。此外还有一组常用的包含关系：$V_1\subseteq V_2$ $\iff$ $V_1\cap V_2=V_1$ $\iff$ $V_1+V_2=V_2$；并且若 $W\subseteq V_1$、$W\subseteq V_2$，则 $W\subseteq V_1\cap V_2$。

\textbf{一个基本等式}：$L(\alpha_1,\dots,\alpha_s)+L(\beta_1,\dots,\beta_t)=L(\alpha_1,\dots,\alpha_s,\beta_1,\dots,\beta_t)$——两个子空间的和，就是把两组向量合并后再张成。
【反哺点】服务考纲〔向量〕“了解子空间的和与交的概念”。求 $V_1+V_2$ 的基，直接合并两组生成向量再求极大无关组即可；求 $V_1\cap V_2$，则令向量同时属于两者，解方程组求公共部分。“$V_1\subseteq V_2$ 的三个等价说法”是证明包含关系的标准工具，比按定义逐向量验证快得多。
\end{fdeep}

% ---------- 维数公式（北大 六.6 定理 7 及推论，印刷 176–177）----------
\begin{fdeep}{维数公式：$\dim(V_1+V_2)=\dim V_1+\dim V_2-\dim(V_1\cap V_2)$}
设 $V_1,V_2$ 是线性空间 $V$ 的两个子空间，则
\fml{\dim(V_1+V_2)=\dim V_1+\dim V_2-\dim(V_1\cap V_2)}

\textbf{证明思路（三段式，值得记住）}：设 $\dim V_1=n_1$，$\dim V_2=n_2$，$\dim(V_1\cap V_2)=m$。第一步，取 $V_1\cap V_2$ 的一组基 $\alpha_1,\dots,\alpha_m$；第二步，由基的扩充定理把它分别扩充成 $V_1$ 的基 $\alpha_1,\dots,\alpha_m,\beta_1,\dots,\beta_{n_1-m}$ 与 $V_2$ 的基 $\alpha_1,\dots,\alpha_m,\gamma_1,\dots,\gamma_{n_2-m}$；第三步，证明把这两组拼起来所得的向量组（共 $n_1+n_2-m$ 个）线性无关，因此它就是 $V_1+V_2$ 的基。证明线性无关时用反证：设线性组合为零并记该向量为 $\alpha$，则 $\alpha\in V_1$ 且 $\alpha\in V_2$，故 $\alpha\in V_1\cap V_2$，可由 $\alpha_1,\dots,\alpha_m$ 表出，代回原式再由两组基各自线性无关即得系数全为零。

\textbf{推论}：若 $\dim V_1+\dim V_2>n=\dim V$，则 $V_1\cap V_2$ 中必含非零向量（因为 $\dim(V_1\cap V_2)=\dim V_1+\dim V_2-\dim(V_1+V_2)$，而 $\dim(V_1+V_2)\le n$）。
【反哺点】服务考纲〔向量〕“了解子空间的交与和”，这是 9 讲几乎不涉及但考研大题常考的公式。三处高频用途：由两组基的维数直接算 $V_1\cap V_2$ 的维数；由“交集为零”反推 $\dim(V_1+V_2)$ 是否等于维数之和（这正是直和）；用推论证明“两个子空间必有公共非零向量”，不必具体解出向量。
\end{fdeep}

% ---------- 直和（北大 六.7 定义 9、10，定理 8、9、10、11，印刷 177–179）----------
\begin{fdeep}{直和的判定：从两个子空间到多个子空间}
若 $V_1+V_2$ 中每个向量 $\alpha$ 的分解式 $\alpha=\alpha_1+\alpha_2$（$\alpha_1\in V_1$，$\alpha_2\in V_2$）都是\textbf{唯一}的，就称这个和为\textbf{直和}，记作 $V_1\oplus V_2$。

\textbf{两个子空间的直和，以下条件等价}（北大 六.7 定理 8、9）：
1）$V_1+V_2$ 是直和；
2）零向量的分解式唯一（即由 $\alpha_1+\alpha_2=\mathbf{0}$，$\alpha_i\in V_i$ 只能推出 $\alpha_1=\alpha_2=\mathbf{0}$）；
3）$V_1\cap V_2=\{\mathbf{0}\}$；
4）$\dim(V_1+V_2)=\dim V_1+\dim V_2$。
其中 2）与 3）的等价最常用：若 $V_1\cap V_2$ 中有非零向量 $\alpha$，则 $\mathbf{0}=\alpha+(-\alpha)$ 就是一个非平凡分解；反之若 $\alpha_1=-\alpha_2\ne\mathbf{0}$，则该向量同时属于两个子空间。
\textbf{推广到多个子空间}（定理 11）：$V_1,V_2,\dots,V_s$ 的和是直和 $\iff$ 零向量表法唯一 $\iff$ 对每个 $i$ 都有 $V_i\cap\sum_{j\ne i}V_j=\{\mathbf{0}\}$ $\iff$ $\dim\sum_{i=1}^{s}V_i=\sum_{i=1}^{s}\dim V_i$。注意多子空间时必须\textbf{逐个}与“其余全部之和”求交，只验两两之交为零是不够的。此外任一子空间 $W$ 都有补子空间 $W'$ 使 $V=W\oplus W'$（定理 10）。
【反哺点】服务考纲〔向量〕“了解子空间的概念”“了解维数”。考纲未直接要求“直和”，但它一旦出现，几乎都落在“$V_1\cap V_2=\{\mathbf{0}\}$”与“维数相加”这两条上：题目给出 $\dim V_1+\dim V_2=\dim V$ 时，要立刻想到“只需再验交为零，即可断言 $V=V_1\oplus V_2$，从而任一向量分解唯一，坐标可以分开算”。这是解“分解唯一性”与“求补空间”两类题的快捷入口。
\end{fdeep}

% ---------- 内积的定义（北大 九.1 定义 1，印刷 245；樊 9.1.1、9.1.3 印刷 580）----------
\begin{fdeep}{内积与欧几里得空间：四条公理与两个典型例子}
设 $V$ 是实数域 $\mathbb{R}$ 上的线性空间。若在 $V$ 上定义了一个二元实函数 $(\alpha,\beta)$，满足
\fml{(\alpha,\beta)=(\beta,\alpha);\qquad (k\alpha,\beta)=k(\alpha,\beta);\qquad (\alpha+\beta,\gamma)=(\alpha,\gamma)+(\beta,\gamma);\qquad (\alpha,\alpha)\ge0}
并且 $(\alpha,\alpha)=0$ 当且仅当 $\alpha=\mathbf{0}$，
就称 $(\alpha,\beta)$ 为 $\alpha$ 与 $\beta$ 的\textbf{内积}，$V$ 称为\textbf{欧几里得空间}（简称欧氏空间）。这四条依次叫\textbf{对称性、齐次性、可加性、正定性}；前三条合起来就是“对第一个变元线性”（再配合对称性即对两个变元都线性），第四条保证长度有意义。

\textbf{两个典型例子}：
（1）$\mathbb{R}^n$ 中 $\alpha=(a_1,\dots,a_n)$，$\beta=(b_1,\dots,b_n)$，规定 $(\alpha,\beta)=a_1b_1+a_2b_2+\cdots+a_nb_n$，称为 $\mathbb{R}^n$ 的\textbf{标准内积}；$n=3$ 时它就是几何向量内积的坐标表达式。
（2）$m\times n$ 实矩阵空间 $M_{m,n}(\mathbb{R})$ 中规定 $(A,B)=\operatorname{tr}(A^{\mathrm T}B)=\sum_{i,j}a_{ij}b_{ij}$，称为 \textbf{Frobenius 内积}。
【反哺点】服务考纲〔向量〕“理解内积的概念”。考纲要求会判断一个二元函数是否为内积，判据就是这四条：四条中任何一条（尤其是正定性）不成立就不是内积。注意正定性里“等号仅当 $\alpha=\mathbf{0}$ 时成立”这一细节，类似 $(A,B)=\operatorname{tr}(AB)$ 或含权为负的“内积”都栽在这一点上。
\end{fdeep}

% ---------- 柯西–施瓦茨不等式（北大 九.1，印刷 246–247；樊 9.1.2 印刷 580）----------
\begin{fdeep}{柯西–施瓦茨不等式：证明套路与等号条件}
对欧氏空间中任意向量 $\alpha,\beta$，有
\fml{\lvert(\alpha,\beta)\rvert\le\lVert\alpha\rVert\,\lVert\beta\rVert}
等号成立当且仅当 $\alpha$ 与 $\beta$ \textbf{线性相关}。（北大称柯西–布尼亚科夫斯基不等式。）

\textbf{证明套路（“造一个非负二次式，看判别式”）}：$\beta=\mathbf{0}$ 时显然。设 $\beta\ne\mathbf{0}$，$t$ 为实变量，作 $\gamma=\alpha+t\beta$。由内积的正定性 $(\gamma,\gamma)\ge0$ 得
\fml{(\alpha,\alpha)+2(\alpha,\beta)t+(\beta,\beta)t^{2}\ge0}
它对一切实数 $t$ 成立。这是开口向上的二次式恒非负，判别式必 $\le0$：$4(\alpha,\beta)^2-4(\alpha,\alpha)(\beta,\beta)\le0$，即得结论。也可直接取 $t=-\dfrac{(\alpha,\beta)}{(\beta,\beta)}$ 代入化简（这正是 $\alpha$ 在 $\beta$ 方向上投影系数的来源）。等号成立 $\iff$ 存在 $t$ 使 $\gamma=\mathbf{0}$ $\iff$ $\alpha,\beta$ 线性相关。

\textbf{两个直接用途}：$\mathbb{R}^n$ 中即 $\lvert\sum a_ib_i\rvert\le\sqrt{\sum a_i^2}\sqrt{\sum b_i^2}$；函数空间中即 $\lvert\int_a^b fg\,\mathrm{d}x\rvert\le\bigl(\int_a^b f^2\,\mathrm{d}x\bigr)^{1/2}\bigl(\int_a^b g^2\,\mathrm{d}x\bigr)^{1/2}$。正因为有它，$\cos\langle\alpha,\beta\rangle=\dfrac{(\alpha,\beta)}{\lVert\alpha\rVert\lVert\beta\rVert}$ 才落在 $[-1,1]$ 内，夹角定义才合法。
【反哺点】服务考纲〔向量〕“了解内积的概念”“掌握施密特方法”。不等式本身常用于求参数范围与证明不等式：把待证式写成内积形式 $(\alpha,\beta)$，再套 $\lvert(\alpha,\beta)\rvert\le\lVert\alpha\rVert\lVert\beta\rVert$。等号条件“线性相关”是很多证明题取等号的依据，也是施密特正交化中“投影系数”分母 $(\beta_i,\beta_i)$ 必须非零（即 $\beta_i\ne\mathbf{0}$）的理论保证。
\end{fdeep}

% ---------- 长度、夹角、正交与勾股（北大 九.1 定义 2、3、4，印刷 246–248）----------
\begin{fdeep}{长度、夹角、正交与勾股定理}
\textbf{长度}：非负实数 $\lVert\alpha\rVert=\sqrt{(\alpha,\alpha)}$ 称为向量 $\alpha$ 的长度（模）。它满足 $\lVert k\alpha\rVert=\lvert k\rvert\,\lVert\alpha\rVert$，且只有零向量的长度为零。长度为 $1$ 的向量称为单位向量；当 $\alpha\ne\mathbf{0}$ 时，$\dfrac{1}{\lVert\alpha\rVert}\alpha$ 是单位向量，这一步叫把 $\alpha$ \textbf{单位化}。

\textbf{夹角}：非零向量 $\alpha,\beta$ 的夹角定义为 $\langle\alpha,\beta\rangle=\arccos\dfrac{(\alpha,\beta)}{\lVert\alpha\rVert\lVert\beta\rVert}$，取值在 $[0,\pi]$。由柯西–施瓦茨不等式，这个余弦值必落在 $[-1,1]$ 内。

\textbf{三角不等式}：$\lVert\alpha+\beta\rVert\le\lVert\alpha\rVert+\lVert\beta\rVert$。证明只需把 $\lVert\alpha+\beta\rVert^2$ 展开，用柯西–施瓦茨不等式控制中间项 $2(\alpha,\beta)\le2\lVert\alpha\rVert\lVert\beta\rVert$，右边恰好是 $(\lVert\alpha\rVert+\lVert\beta\rVert)^2$。

\textbf{正交}：若 $(\alpha,\beta)=0$，就称 $\alpha$ 与 $\beta$ \textbf{正交}（互相垂直），记作 $\alpha\perp\beta$。零向量与任何向量（包括它自己）都正交，所以讨论正交向量组时通常要求\textbf{非零}。两个非零向量正交 $\iff$ 它们的夹角为 $\dfrac{\pi}{2}$。

\textbf{勾股定理}：当 $\alpha\perp\beta$ 时 $\lVert\alpha+\beta\rVert^2=\lVert\alpha\rVert^2+\lVert\beta\rVert^2$；更一般地，若 $\alpha_1,\dots,\alpha_m$ 两两正交，则 $\lVert\alpha_1+\cdots+\alpha_m\rVert^2=\lVert\alpha_1\rVert^2+\cdots+\lVert\alpha_m\rVert^2$。
【反哺点】服务考纲〔向量〕“理解内积的概念，掌握用内积判断向量的正交性”。正交的判据就是内积为零，而内积在标准正交基下化简为坐标分量乘积之和，所以“判断正交”常常转化为“算一个点积”。勾股定理的推广形式是“正交组”题目的万能展开式，比逐项算内积快得多。
\end{fdeep}

% ---------- 度量矩阵：内积的坐标表达式与换基合同（北大 九.1，印刷 248–249）----------
\begin{fdeep}{度量矩阵：把内积写成矩阵，换基时变合同}
设 $\varepsilon_1,\dots,\varepsilon_n$ 是 $n$ 维欧氏空间 $V$ 的一组基，$a_{ij}=(\varepsilon_i,\varepsilon_j)$，则 $A=(a_{ij})$ 称为这组基的\textbf{度量矩阵}（它就是向量组 $\varepsilon_1,\dots,\varepsilon_n$ 的 Gram 矩阵；其定义、秩公式与正定性的完整讨论见本章 Gram 矩阵深化块，此处不重复）。

有了它，内积可以\textbf{按坐标直接算}：若 $\alpha=\sum_i x_i\varepsilon_i$，$\beta=\sum_j y_j\varepsilon_j$，则
\fml{(\alpha,\beta)=\sum_{i,j}a_{ij}x_iy_j=X^{\mathrm T}AY}
其中 $X=(x_1,\dots,x_n)^{\mathrm T}$、$Y=(y_1,\dots,y_n)^{\mathrm T}$ 分别是 $\alpha,\beta$ 的坐标。把内积展开，这正是一个以 $A$ 为矩阵的二次型：$(\alpha,\alpha)=X^{\mathrm T}AX$。所以\textbf{度量矩阵完全确定了内积}；这也说明“$A$ 正定”与“长度的平方恒非负、且仅零向量为零”是同一句话。

\textbf{换基时的变化}：若从 $\varepsilon_1,\dots,\varepsilon_n$ 到 $\eta_1,\dots,\eta_n$ 的过渡矩阵为 $C$，则新基的度量矩阵 $B$ 满足
\fml{B=C^{\mathrm T}AC}
即不同基下的度量矩阵彼此\textbf{合同}。请对照二次型的换元 $X=CY$：此时矩阵同样变成 $C^{\mathrm T}AC$，两处是同一个代数事实。标准正交基之所以好算，正因为它的度量矩阵是单位矩阵 $E$，$(\alpha,\beta)=X^{\mathrm T}Y$ 退化成普通点积。
【反哺点】服务考纲〔向量〕“理解内积的概念”，并与〔二次型〕“了解合同变换”直接呼应。三个考场用途：题目给出基的内积表（即度量矩阵）时，一切内积计算化为 $X^{\mathrm T}AY$，不必逐个展开；“合同”与“相似”的区分中，度量矩阵换基正是合同的教科书例子；由 $B=C^{\mathrm T}AC$ 知 $|B|=|C|^{2}|A|$，可判断换基前后度量矩阵行列式的符号与是否为零都不变。
\end{fdeep}

% ---------- 标准正交基（北大 九.2 定义 5、6，印刷 249–250；樊 9.1.4 印刷 580）----------
\begin{fdeep}{标准正交基：正交组必线性无关，坐标与内积同时变简单}
\textbf{定义}：欧氏空间 $V$ 中一组\textbf{非零}向量若两两正交，称为\textbf{正交向量组}；由 $n$ 个向量组成的正交向量组称为\textbf{正交基}；由单位向量组成的正交基称为\textbf{标准正交基}（规范正交基），即满足 $(\eta_i,\eta_j)=\delta_{ij}$。

\textbf{正交向量组必线性无关}。证明只有两步：设 $k_1\alpha_1+\cdots+k_m\alpha_m=\mathbf{0}$，用 $\alpha_i$ 与等式两边作内积，由两两正交得 $k_i(\alpha_i,\alpha_i)=0$；又 $\alpha_i\ne\mathbf{0}$ 故 $(\alpha_i,\alpha_i)>0$，于是 $k_i=0$。由此立即得到：$n$ 维欧氏空间中两两正交的非零向量\textbf{不可能超过 $n$ 个}。

\textbf{标准正交基的三条好处}：
1）一组基是标准正交基 $\iff$ 它的度量矩阵是单位矩阵 $E$（因为 $a_{ij}=(\eta_i,\eta_j)=\delta_{ij}$）。标准正交基一定存在。
2）坐标可以“用内积直接读出”：若 $\alpha=x_1\eta_1+\cdots+x_n\eta_n$，则 $x_i=(\eta_i,\alpha)$，即 $\alpha=(\eta_1,\alpha)\eta_1+\cdots+(\eta_n,\alpha)\eta_n$。
3）内积化为坐标的普通点积：$(\alpha,\beta)=x_1y_1+x_2y_2+\cdots+x_ny_n=X^{\mathrm T}Y$。这就是“在标准正交基下，欧氏空间的几何就是 $\mathbb{R}^n$ 的几何”的原因。
【反哺点】服务考纲〔向量〕“了解规范正交基的概念及其性质”。选择题常考“正交向量组是否线性无关”（是）、“$n$ 维空间中能否有 $n+1$ 个两两正交的非零向量”（不能）。第 2 条把“求坐标”变成“算内积”，在已知标准正交基的题目中比解方程组快得多；第 3 条则是实对称矩阵正交相似对角化后内积仍好算的依据。
\end{fdeep}

% ---------- 正交组的扩充与过渡矩阵（北大 九.2 定理 1 印刷 250；印刷 252）----------
\begin{fdeep}{正交组的扩充：定理 1 的构造法本身就是施密特方法}
\textbf{定理 1}：$n$ 维欧氏空间中任一正交向量组都能扩充成一组正交基。

证明用的是对 $n-m$ 的归纳，但\textbf{归纳步给出的就是构造方法}：设已有正交组 $\alpha_1,\dots,\alpha_m$，取一个不能被它们线性表出的向量 $\beta$，令
\fml{\alpha_{m+1}=\beta-k_1\alpha_1-k_2\alpha_2-\cdots-k_m\alpha_m,\qquad k_i=\frac{(\beta,\alpha_i)}{(\alpha_i,\alpha_i)}}
这里的 $k_i$ 正是 $\beta$ 在 $\alpha_i$ 方向上的\textbf{投影系数}，故 $(\alpha_{m+1},\alpha_i)=(\beta,\alpha_i)-k_i(\alpha_i,\alpha_i)=0$，即 $\alpha_{m+1}$ 与每个 $\alpha_i$ 都正交。又因 $\beta$ 不能被 $\alpha_1,\dots,\alpha_m$ 表出，故 $\alpha_{m+1}\ne\mathbf{0}$。如此逐个扩充，最后单位化即得标准正交基。把“已知正交组”换成“只单位化了前 $m$ 个”的情形，同一套式子逐次做下去就是\textbf{施密特正交化}（其完整叙述见主控撰写的对应拓展块）。

\textbf{一个常被忽略的等价说法}：两组标准正交基之间的\textbf{过渡矩阵必是正交矩阵}。设 $\eta_1,\dots,\eta_n$ 与 $\varepsilon_1,\dots,\varepsilon_n$ 都是标准正交基，$(\eta_1,\dots,\eta_n)=(\varepsilon_1,\dots,\varepsilon_n)A$，则 $A$ 的第 $j$ 列是 $\eta_j$ 在标准正交基 $\varepsilon_i$ 下的坐标，由 $(\eta_i,\eta_j)=\delta_{ij}$ 代入内积的坐标表达式即得 $A^{\mathrm T}A=E$。
【反哺点】服务考纲〔向量〕“掌握线性无关向量组正交规范化的施密特方法”与“了解正交矩阵”。扩充定理的公式 $\alpha_{m+1}=\beta-\sum\frac{(\beta,\alpha_i)}{(\alpha_i,\alpha_i)}\alpha_i$ 就是施密特方法的单步，理解了它，正交化公式就不必硬背（“减掉已有方向上的投影”）。而“标准正交基到标准正交基的过渡矩阵是正交矩阵”把第 6 章的基变换与第 7、8 章的正交对角化串成一条线：实对称矩阵正交对角化时拼出的 $Q$ 正是两组标准正交基之间的过渡矩阵。
\end{fdeep}


% ===========================================================================
%  【主控裁决 · 第 6 章 deep_ch06 取舍】（依据 AGENTS.md §2 决定二）
%
%  第 6 讲骨架仅 5 页（印刷 64–71），深化料须精选，否则喧宾夺主。
%
%  【保留 14 块】
%   ✓ 块1　{线性空间：从具体对象到八条公理
%   ✓ 块2　{子空间的判别：两条封闭性就够
%   ✓ 块3　{生成子空间 $L(\alpha_1,\dots,\alpha_s)$：极大无关组就是它的基
%   ✓ 块5　{线性相关性的判定网：五条入口殊途同归
%   ✓ 块6　{极大线性无关组与秩：它和“基”是同一个东西
%   ✓ 块7　{向量组等价与矩阵等价：两个“等价”不是一回事
%   ✓ 块10　{维数公式：$\dim(V_1+V_2)=\dim V_1+\dim V_2-\dim(V_1\cap V_2)$
%   ✓ 块11　{直和的判定：从两个子空间到多个子空间
%   ✓ 块12　{内积与欧几里得空间：四条公理与两个典型例子
%   ✓ 块13　{柯西–施瓦茨不等式：证明套路与等号条件
%   ✓ 块14　{长度、夹角、正交与勾股定理
%   ✓ 块15　{度量矩阵：把内积写成矩阵，换基时变合同
%   ✓ 块16　{标准正交基：正交组必线性无关，坐标与内积同时变简单
%   ✓ 块17　{正交组的扩充：定理 1 的构造法本身就是施密特方法
%
%  【删除 3 块及理由】
%   ✗ 块4　{线性相关与线性无关：两种定义与“部分、整体”性质
%   ✗ 块8　{有限维空间的基与维数：唯一性、判定与扩充
%   ✗ 块9　{子空间的交与和：两种运算与包含关系
%
%   · 块4「线性相关与线性无关：两种定义与部分/整体性质」→ 层次偏低的正文级铺垫，
%       且 9 讲第 6 讲正文已覆盖线性相关性的判定。
%   · 块8「有限维空间的基与维数：唯一性、判定与扩充」→ 与 9 讲「四、向量空间（仅数学一）」
%       已给的基/维数/坐标唯一表示重复。
%   · 块9「子空间的交与和」→ 与块10「维数公式」内容连续，后者已含交与和的运算与包含关系。
%
%  【保留重点（对齐考纲三.5/6/7/8）】
%   向量空间八条公理、子空间判别、生成子空间 L（极大无关组=基）、向量组等价 vs 矩阵等价、
%   **维数公式**、**直和判定**、内积四公理、**柯西–施瓦茨不等式**、长度/夹角/正交、
%   **标准正交基**、**正交组扩充 = 施密特正交化**、度量矩阵。
%   其中内积、柯西不等式、标准正交基、施密特正交化四项正是 9 讲完全缺失的考纲内容。
% ===========================================================================
